% !TEX program = xelatex
\documentclass[12pt,a4paper,fontset=none]{ctexart}
\setCJKmainfont{SimSun}\setCJKsansfont{SimSun}\setCJKmonofont{SimSun}
\setCJKfamilyfont{zhsong}{SimSun}\setCJKfamilyfont{zhhei}{SimSun}
\setCJKfamilyfont{zhkai}{SimSun}

\usepackage[top=2.2cm,bottom=2.2cm,left=2.3cm,right=2.3cm]{geometry}
\usepackage{amsmath,amssymb,booktabs,multirow,array,graphicx,float,hyperref,enumitem,xcolor,colortbl}

\begin{document}

\title{\heiti\zihao{2} 高铁快运基地选址与规模多维协同优化\\[4pt]
       \large 全部问题综合分析报告}
\author{\kaishu\zihao{4} 胡鹏禹}
\date{\zihao{4} 2026年6月7日}
\maketitle
\thispagestyle{empty}

\tableofcontents
\newpage

% ============================================================
\section{问题一：货流最短路与全有全无分配}
% ============================================================

\subsection{问题描述}
假设全网快运货流全部采用货运动车组模式运输。计算每支OD在34节点路网中的最短路径，将1908.4吨/日需求全部沿最短路径分配，检验各区段是否超过线路通过能力（8列/日 $\times$ 85吨/列 = 680吨/日）。

\subsection{路网结构}
区域路网包含34个节点（15个主站+19个中间路由节点）和44条无向边。

\begin{table}[H]\centering\small
\caption{主站连接度（缩略图）}\label{tab:degree}
\begin{tabular}{lclcl}
\toprule
站点 & 相邻主站数 & 站点 & 相邻主站数 \\
\midrule
南京 NJ & 5（BB,HF,HZ,NC,NT） & 武汉 WH & 5（CS,HF,NC,YC,ZZ）\\
合肥 HF & 4（BB,NC,NJ,WH） & 郑州 ZZ & 3（CQ,WH,XZ）\\
重庆 CQ & 4（CD,GY,YC,ZZ） & 南昌 NC & 4（CS,HF,NJ,WH）\\
长沙 CS & 3（GY,NC,WH） & 成都 CD & 1（CQ）\\
\bottomrule
\end{tabular}
\end{table}

成都仅通过重庆与全网连通，是路网中拓扑最脆弱的主站。

\subsection{分配结果}

\begin{table}[H]\centering\small
\caption{全有全无分配核心指标}\label{tab:p1}
\begin{tabular}{lr}
\toprule
指标 & 数值 \\
\midrule
总需求处理量 & 1908.4 吨/日 \\
总周转量 & 1,871,085 吨$\cdot$公里/日 \\
平均运距 & 980 km \\
\rowcolor{green!5}
最大断面流量 & 535.0 吨/日（合肥$\to$六安，饱和度78.7\%） \\
超限区段数 & \textbf{0}（全部在680吨限制内） \\
高负荷区段数（80\%+） & \textbf{0} \\
\bottomrule
\end{tabular}
\end{table}

\begin{table}[H]\centering\small
\caption{OD距离分布}\label{tab:dist}
\begin{tabular}{cccc}
\toprule
距离段 & OD对数 & 需求占比 & 累计需求(吨) \\
\midrule
0--300 km & 8 & 4.1\% & 77.4 \\
300--500 km & 30 & 16.5\% & 315.5 \\
500--800 km & 42 & 20.1\% & 383.2 \\
800--1000 km & 34 & 18.8\% & 358.2 \\
1000--1500 km & 44 & 22.9\% & 436.5 \\
1500--2000 km & 30 & 10.7\% & 204.9 \\
2000+ km & 2 & 7.0\% & 132.6 \\
\bottomrule
\end{tabular}
\end{table}

\subsection{结论}
\textbf{现有高铁路网的线路通过能力十分充裕}，没有任何区段超限，最大饱和度仅78.7\%。路网瓶颈不在线路而在车站处理能力和捎带容量约束。

\clearpage

% ============================================================
\section{问题二：最小费用最大流分配}
% ============================================================

\subsection{模型设计}
在问题一"无约束"基础上，引入两类硬约束：(1)线路通过能力（680吨/日/区段）；(2)车站捎带能力（每站发出/到达各$\le48$吨/日）。使用Gurobi求解"最大化总承运量→最小化总费用"的二阶段LP模型。

\subsection{求解结果}

\begin{table}[H]\centering\small
\caption{问题二核心指标}\label{tab:p2}
\begin{tabular}{lr}
\toprule
指标 & 数值 \\
\midrule
总需求 & 1908.4 吨/日 \\
\rowcolor{green!5}
最大满足量 & \textbf{1908.4 吨/日（100.0\%）} \\
未满足量 & 0.0 吨/日 \\
最小总费用 & 6,453,039 元/日（$\approx$645万/日） \\
\midrule
EMU承运量 & 1212.6 吨/日（63.5\%） \\
捎带承运量 & 695.8 吨/日（36.5\%） \\
\midrule
平均吨成本 & 3,381 元/吨 \\
平均吨收入 & 3,801 元/吨 \\
\rowcolor{red!5}
\textbf{吨均毛利} & \textbf{420 元/吨} \\
\bottomrule
\end{tabular}
\end{table}

\subsection{与问题三的对比——利润从何而来？}

问题二吨均毛利仅420元，但问题三吨均毛利高达2748元。差异根源：

\begin{enumerate}
    \item \textbf{问题二满足100\%需求}，其中大量长尾OD（$<5$吨/日、$>1500$km）的运输收入低于EMU运行成本，形成"倒挂"（吨均毛利被拉低）；
    \item \textbf{问题三通过基地选址约束}，模型战略性放弃58.8\%的低收益货流，仅承运高利润OD（吨均毛利跃升至2748元）；
    \item 这验证了"利润最大化$\neq$需求最大化"的运筹优化核心理念。
\end{enumerate}

\section{问题三：基地选址与规模优化（核心MIP模型）}

\subsection{选址决策}
12个备选站中精准选择\textbf{7个内陆枢纽}，全部"改建小规模"。

\begin{table}[H]\centering\small
\caption{选址与规模决策}\label{tab:loc}
\begin{tabular}{ccl|ccl}
\toprule
站点 & 名称 & 决策 & 站点 & 名称 & 决策 \\
\midrule
WH & 武汉 & \cellcolor{green!15}改建小规模 & NJ & 南京 & \cellcolor{red!10}不建 \\
CS & 长沙 & \cellcolor{green!15}改建小规模 & HZ & 杭州 & \cellcolor{red!10}不建 \\
NC & 南昌 & \cellcolor{green!15}改建小规模 & SH & 上海 & \cellcolor{red!10}不建 \\
ZZ & 郑州 & \cellcolor{green!15}改建小规模 & HF & 合肥 & \cellcolor{red!10}不建 \\
CQ & 重庆 & \cellcolor{green!15}改建小规模 & XZ & 徐州 & \cellcolor{red!10}不建 \\
GY & 贵阳 & \cellcolor{green!15}改建小规模 & & & \\
CD & 成都 & \cellcolor{green!15}改建小规模 & & & \\
\bottomrule
\end{tabular}
\end{table}

\textbf{格局特征}：建成基地全部位于中西部（武汉--长沙--贵阳--成都--重庆"品"字形走廊），东部沿海（上海、杭州、南京）全部放弃。

\subsection{经济指标}

\begin{table}[H]\centering\small
\caption{日均经济指标}\label{tab:econ}
\begin{tabular}{lrr}
\toprule
指标 & 金额（元/日） & 占比 \\
\midrule
运输总收入 & 2,808,289 & 100\% \\
\quad EMU收入 & 495,156 & 17.6\% \\
\quad 捎带收入 & 2,313,133 & 82.4\% \\
\midrule
总成本 & 649,349 & 23.1\% \\
\rowcolor{green!5}
\textbf{日均净利润} & \textbf{2,158,940} & \textbf{76.9\%} \\
\bottomrule
\end{tabular}
\end{table}

\subsection{货流分担格局}

\begin{table}[H]\centering\small
\caption{货流模式分担}\label{tab:flow}
\begin{tabular}{lrr}
\toprule
模式 & 吨/日 & 占比 \\
\midrule
总需求 & 1908.4 & 100\% \\
实际满足 & 785.6 & 41.2\% \\
\quad EMU承运 & 95.8 & 12.2\% \\
\quad 捎带承运 & 689.8 & 87.8\% \\
战略性放弃 & 1122.8 & 58.8\% \\
\bottomrule
\end{tabular}
\end{table}

\subsection{三大核心发现}

\begin{enumerate}
\item \textbf{捎带是利润引擎}。捎带边际成本仅120元/吨（无运行费、无折旧），87.8\%的货流走捎带，贡献82.4\%收入。EMU全成本$2000-3500$元/吨，仅在内陆枢纽间因捎带能力饱和时启用；
\item \textbf{东部零基地是理性选择}。上海、杭州、南京等东部枢纽的捎带频次极高，48吨/日/站即可覆盖主要货流。建库的边际收益远低于边际成本；
\item \textbf{改建小规模碾压新建}。从"改建小规模"到"新建小规模"，日均摊销成本增10倍（1.08万$\to$11.1万），吞吐能力仅增4倍（1列$\to$4列）。在当前需求密度下，新建方案的投资回报率为负。
\end{enumerate}

\clearpage

\section{问题四：列车开行方案设计}

基于问题三MIP模型导出的边流量和捎带流量，进行运输服务拆解。

\subsection{货运动车组班列方案}

\begin{table}[H]\centering\small
\caption{EMU骨干班列}\label{tab:emu}
\begin{tabular}{lcc}
\toprule
运行路由 & 距离(km) & 开行频次 \\
\midrule
长沙$\to$贵阳$\to$成都$\to$重庆$\to$郑州 & 2715 & 双向各1列/日 \\
长沙$\to$贵阳$\to$成都 & 1354 & 双向各1列/日 \\
长沙$\to$贵阳$\to$重庆$\to$郑州 & 2120 & 双向各1列/日 \\
武汉$\leftrightarrow$长沙直达 & 362 & 双向各1列/日 \\
\bottomrule
\end{tabular}
\end{table}

共17条EMU运行线，17列/日，承运125.3吨（内陆走廊）。

\subsection{捎带开行方案}

共119条捎带OD线，151列/日，承运689.8吨。主力OD为长距离跨区（如杭州$\to$成都1900km, 17.3吨/日）。

\subsection{能力校核}
区间EMU列车最大4列/日（限制8列）$\checkmark$；长沙站等效占用14列/日（捎带+EMU）为全网最高。

\section{问题五：大周期盈亏平衡与财务评估}

\subsection{投资概况}

\begin{table}[H]\centering\small
\caption{投资与回报}\label{tab:invest}
\begin{tabular}{lr}
\toprule
指标 & 数值 \\
\midrule
建设基地 & 7个改建小规模 \\
单站投资 & 0.768亿元 \\
总投资 & \textbf{5.376亿元} \\
年固定人工 & 84万元 \\
\midrule
年净利润 & \textbf{7.91亿元} \\
投资回收期 & \textbf{8.2个月} \\
ROI & \textbf{147\%/年} \\
20年NPV（折现率5\%） & \textbf{93.1亿元} \\
IRR & \textbf{147\%} \\
盈亏平衡日货量 & 27吨（占实际3.4\%） \\
安全边际 & \textbf{96.6\%} \\
\bottomrule
\end{tabular}
\end{table}

\subsection{敏感性分析}
\begin{itemize}
    \item 需求波动-40\%：年净利降至4.63亿，回收期1.16年——\textbf{仍强健}
    \item 运价波动-30\%：年净利降至7.36亿——\textbf{极低风险}
    \item 建设成本+30\%：回收期仅延至10.7月——\textbf{投资安全}
\end{itemize}

\section{问题六：多式联运市场竞争分析}

\subsection{纯运输成本对比（200吨货物）}

\begin{table}[H]\centering\small
\caption{三种模式成本对比}\label{tab:modes}
\begin{tabular}{cccc}
\toprule
距离 & 高铁快运（万元） & 航空货运（万元） & 公路卡车（万元） \\
\midrule
300 km & 28.4 & 114.0 & \cellcolor{green!10}\textbf{5.5} \\
800 km & 45.9 & 154.0 & \cellcolor{green!10}\textbf{10.8} \\
1500 km & 70.3 & 210.0 & \cellcolor{green!10}\textbf{18.1} \\
2000 km & 87.8 & 250.0 & \cellcolor{green!10}\textbf{23.4} \\
\bottomrule
\end{tabular}
\end{table}

公路在纯运输成本上全距离占优。高铁固定成本（50000元/列）远高于公路（2000元/车）。

\subsection{综合竞争力}

\begin{table}[H]\centering\small
\caption{六维竞争力评分}\label{tab:radar}
\begin{tabular}{lccc}
\toprule
维度 & 高铁 & 航空 & 公路 \\
\midrule
运输时效 & 8 & 10 & 2 \\
运价竞争力 & 6 & 2 & 10 \\
准点可靠性 & \textbf{10} & 5 & 6 \\
装载能力 & 7 & 3 & 8 \\
绿色低碳 & \textbf{9} & 2 & 6 \\
网络覆盖 & 5 & 4 & \textbf{10} \\
\midrule
综合平均 & \textbf{7.5} & 4.3 & 7.0 \\
\bottomrule
\end{tabular}
\end{table}

\subsection{结论}
高铁的竞争优势不在价格而在\textbf{准点率（99\% vs 航空75\%）和低碳（碳排放为航空1/30）}。推荐定位500--1500km高附加值时效敏感市场。

% ============================================================
\section{综合结论与建议}
% ============================================================

\subsection{六道问题的逻辑递进关系}

\textbf{问题一}（无约束分配）$\to$ 确认路网线路能力充裕，瓶颈不在线路。

\textbf{问题二}（有约束MCMF）$\to$ 100\%需求可满足，但吨利仅420元——"满额承运"并不等于"利润最大"。

\textbf{问题三}（MIP选址）$\to$ 核心模型。7个内陆改建小规模基地 + 捎带优先策略，净利216万/日，验证"战略性放弃=利润最大化"。

\textbf{问题四}（开行方案）$\to$ 将问题三的流分配落地为可执行列车计划。

\textbf{问题五}（财务评估）$\to$ 回收期8.2月，NPV 93亿，项目极强健。

\textbf{问题六}（竞争分析）$\to$ 高铁不拼价格，拼时效+可靠性+低碳。

\subsection{政策建议}

\begin{enumerate}[leftmargin=*]
    \item \textbf{短期（1--2年）}：聚焦提升捎带能力——增加确认车班次。这是ROI最高的投资方向；
    \item \textbf{中期（3--5年）}：在需求增长至2500+吨/日后，内陆基地逐步扩容至"新建小规模"；
    \item \textbf{长期}：警惕公路时效提升和航空降价对高铁500--1500km优势带的侵蚀。
\end{enumerate}

\end{document}
